\documentclass[../main.tex]{subfiles}

\begin{document}

%% ======================================================================
%%  Section 8 – Independent HDF5 Implementations
%% ======================================================================
\section{Independent HDF5 Implementations}
\label{sec:indie}

\subsection{Overview}

\begin{description}
    \item[PureHDF] Managed .NET read/write without native binaries.
    \item[jHDF] JVM HDF5 reading.
    \item[netCDF-Java] JVM netCDF/HDF-EOS/HDF5 scientific-data reading.
    \item[pyfive] Pure Python read-only inspection.
    \item[jsfive] Browser-side direct HDF5 reading.
    \item[scigolib/hdf5] Pure Go read/write.
    \item[hdf5-pure and rustyhdf5-format/rustyhdf5] Pure Rust read/write.
    \item[async-hdf5 or h5coro] Cloud object-store chunk mapping.
    \item[JLD2.jl] Julia-native persistence in HDF5-compatible files.
    \item[hdf5-pickles / H5Lens] Machine-readable format definitions with an
          independent reader, validator, and preflight oracle. Distinct from
          the others in that its interpretation of the format is
          \emph{inspectable as data} rather than inferable from behaviour.
\end{description}

\paragraph{Vocabulary.}
This chapter uses the issue families defined by the SSP SIG's HDF5 Independent
Implementation (H5II) model~\cite{hdf5-indie-model}: \textbf{I1} parser/spec
divergence, \textbf{I2} fail-open compatibility behaviour, \textbf{I3} writer
correctness and interoperability failure, \textbf{I4} codec, filter, or
external-reference boundary failure, \textbf{I5} numeric, structural, or
resource amplification, \textbf{I6} concurrency, durability, or lifecycle
mismatch, \textbf{I7} metadata and artifact exposure, \textbf{I8} dependency,
provenance, or compatibility-claim drift, and \textbf{I9} trust-boundary
confusion. As elsewhere, these labels classify mechanisms; they are not risk
ratings, and they do not carry the SSP SIG's numeric scores into
Section~\ref{sec:risk-rubric}.

The model's central observation frames this chapter: an independent
implementation removes a dependency on the reference library but not the hard
problem, because the trusted computing base becomes the implementation's own
parser, validator, writer, and support libraries. Its working assumption that
\emph{partial support is normal, but unsupported or unknown features must fail
closed rather than being silently ignored or guessed} is the specific standard
against which the two 2026 cases below are best read.

\paragraph{Why the count matters.}
The number of independent implementations is not a curiosity; it changes the
character of specification risk. With a single implementation, an ambiguity in
the format is invisible, because that implementation's behaviour is
operationally indistinguishable from the format itself --- there is nothing for
it to disagree with. With five or nine, \textbf{every ambiguity is resolved
independently once per implementation}, and each unresolved ambiguity becomes a
latent divergence that surfaces only when files cross implementation
boundaries. That is precisely when the data is least recoverable and the
diagnosis hardest, because the failure appears as a reader defect rather than
as the specification defect it is.

This is the mechanism behind \code{FMT-SAF-001}. Both documented instances ---
issues \#6409 and \#6616 --- were found by an independent implementation
exercising a construct and diverging from the reference library, not by reading
the specification. Independent implementations are therefore currently the
ecosystem's only conformance test, and what they are detecting is not confined
to their own defects: they are finding errors in the reference library and in
the specification. That is a strength of the ecosystem and an indictment of the
normative artifact at the same time.

Independent readers and writers are valuable checks on both the prose
specification and the reference implementation. They also expose a three-way
distinction essential to interoperability analysis: a file can be accepted by
one reader, conform to or violate the normative format independently of that
acceptance, and remain safe or unsafe for a particular deployment. Reader
acceptance alone is not conformance evidence.

\subsection{Two Distinct 2026 Chunk-Layout Cases}

The two cases below produced superficially similar upgrade failures, but they
must not be combined. One involved a reader imposing a stronger condition than
the format required; the other involved contradictory writer output that older
readers had tolerated.

\subsubsection{Sufficient but Nonminimal Encoding}

HDF5 issue \#6409 concerns version-4 chunk-layout dimension widths. The format
required an encoded width sufficient to hold the dimension value, but did not
require the smallest possible width. The exact-width check present in HDF5
2.0.0, 2.1.0, and 2.1.1 could therefore reject a specification-permitted,
nonminimal encoding produced by an independent writer. Stated plainly, and this
is the part the interoperability framing tends to soften: \textbf{the reference
library was non-conforming}, the specification was correct, and it took an
independent implementation to establish that. The merged HDF5-side
change accepts any sufficient width and was tracked against 2.2.0
\cite{hdf5-gh-issue-6409,hdf5-gh-pr-6432}. PureHDF separately changed new
output to match the reference reader's expectation
\cite{purehdf-gh-issue-168,purehdf-pr-160} --- a rational local decision with a
corrosive global effect, since converging on the reference implementation
rather than on the specification is how a specification quietly stops being
normative. \code{FMT-SAF-001} records that dynamic.
Formal finding \code{IND-SAF-001} records the historical reader-side
interoperability risk; issue closure is not, by itself, verification of every
release artifact or backport. In H5II terms this is \textbf{I1}, parser/spec
divergence --- and note that it is divergence by the \emph{reference} reader,
which is the direction the model's framing makes it easy to miss.

\subsubsection{Contradictory PureHDF Metadata}

HDF5 issue \#6534 began as a suspected HDF5 2.0 regression because the same
production-derived file opened with HDF5 1.10.10 and 1.14.6 but failed with
HDF5 2.0. Maintainer byte-level analysis identified a different condition in
the affected PureHDF 2.1.2 output: float64 datatype metadata declared an
eight-byte element while the version-4 chunk-layout message recorded a
four-byte stored element. The file was internally inconsistent; the older
readers had been permissive, and the HDF5 2.0 rejection was correct
\cite{hdf5-gh-issue-6534,purehdf-gh-issue-164}.

The reporter described a continuously written production acquisition pipeline
and a reader-version pin as the contemporaneous workaround. A technical metadata-
repair path was proposed and older readers remained a workaround, although the
reporter said rewriting the files was not operationally feasible in that
pipeline. The pin restored availability but was neither a conformance repair
nor the report's Immediate urgency category. The evidence therefore supports
bounded operational and archival
harm, not irreversible data loss. Formal finding \code{IND-SAF-002} records
this writer-side condition. Existing affected files remain a migration and
inventory question even after corrected writer output is deployed. In H5II
terms it is \textbf{I3}, writer correctness and interoperability failure, and
the older readers' acceptance of the file is \textbf{I2}, fail-open
compatibility behaviour --- the pair is the model's canonical chain, where a
writer emits bytes that are only self-consistent to itself and a permissive
reader conceals it until a stricter one does not.


\subsection{When the Specification Itself Is the Defect}

The two chunk-layout cases above are divergences with a determinable answer:
once examined, it was clear which party was wrong. A third case shows that this
is not guaranteed.

HDF5 issue \#6616 records that the file-format specification documents the
group info message as \emph{optional}, while the C library raises a ``message
type not found'' error when it is absent and consequently cannot modify a file
written without one. The reporter --- again an independent-implementation
maintainer, again working from a PureHDF change --- sets out the two available
resolutions: the library should fall back to default group-info values, or the
specification should be corrected to state that the message is required. At the
evidence cutoff the issue is open, labelled a documentation bug, assigned, and
milestoned, with no ruling recorded~\cite{hdf5-gh-issue-6616}.

Until that ruling exists, \textbf{it is not determinable whether a file without
a group info message conforms}. This matters beyond the individual case for two
reasons. First, it is a limitation of this report's own method: the conformance
analysis in Section~\ref{sec:assessment-methodology} distinguishes the
normative requirement from writer behaviour and reader validation, and here the
normative requirement is exactly what is in dispute. Second, it shows that the
missing control is not only a better artifact but a \emph{procedure} --- who
rules, on what timeline, and where the ruling is recorded. A machine-readable
specification would have forced the question earlier by making the optionality
explicit and testable; it would not, by itself, have answered it.

The condition is registered as \code{FMT-SAF-001} in
Appendix~\ref{app:findings-register} and treated by \code{FMT-S1} below.

\subsection{Audit Findings and Mitigation Planning}

The main interoperability challenge across direct-format implementations is not
opening simple files; it is coverage of the long tail: dense groups, fractal
heaps, shared object-header messages, chunk-index variants, references, VLEN,
compound and nested datatypes, virtual datasets, external storage, SWMR flags,
user blocks, non-core filters, unusual library-version bounds, and files
produced by different HDF5 releases. This architectural observation is not
assigned a risk tier without a concrete affected workflow. The H5II model's
\code{II-\#\#\#} review register is the natural instrument for working through
that long tail feature by feature, and \code{IND-S1} below supplies the
fixtures it would need. The two bounded
formal findings are \code{IND-SAF-001} and \code{IND-SAF-002} in
Appendix~\ref{app:findings-register}; the condition that produces them is
registered separately as \code{FMT-SAF-001}.

\paragraph{Priority reconciliation.}\mbox{}\\
\code{FMT-SAF-001} is routed Near-term, and its Near-term obligation is narrow:
rule on the disputes already identified, record each ruling, and change
whichever artifact was wrong. That is a small, bounded action with a named
owner in at least one case, and it is not contingent on \code{FMT-S1}.
\code{IND-SAF-001} and \code{IND-SAF-002} are Moderate and follow the Planned
route; neither is ordered
ahead of the other by issue chronology or recommendation numbering. Phase~2
therefore verifies the exact supported reader artifacts, preserves positive and
negative fixtures, inventories the affected legacy PureHDF population, and
provides a deterministic repair, migration, or compatibility disposition.
Invariant and sustainment work continues in Phases~3--4. A production or
archival population discovered to be actively interrupted may be expedited,
but that scoped decision does not revise the historical inherent ratings.

The \code{FMT-S} and \code{IND-S} identifiers below identify recommendations
only; their numbers do not express risk, urgency, effort, dependency, or
implementation order.

\subsubsection*{Recommendation FMT-S1: Make the specification machine-readable and adjudicable}

Two separate deliverables are conflated whenever this is discussed, and they
fail in different ways if either is attempted alone.

\paragraph{The artifact.}
Maintain the normative file-format description in a machine-readable form from
which validators and fixtures can be generated, rather than only as prose. The
argument is not tooling convenience. \textbf{Prose can be silently ambiguous;
an executable description cannot.} Writing ``the group info message is
optional'' in a form that must generate a valid file and validate one forces
the question that issue \#6616 is currently suspended on --- optional to write,
optional to read, or optional only when some other condition holds. The
existing \code{hdf5-pickles} work is the obvious foundation, and it is already
an independent implementation in its own right rather than only a fixture
generator~\cite{hdf5-pickles}.

Two constraints on that foundation. It is GPLv3 while HDF5 is BSD-style, so
the relationship must be consumption of identifiers, boundaries, and fixtures
as data rather than incorporation of source; see
Section~\ref{sec:roadmap-licensing}. And a machine-readable description is
normative only if the project says it is --- otherwise it becomes a fourth
artifact to disagree with, alongside the prose, the library, and the
implementations.

\paragraph{The procedure.}
Adopt a standing rule for specification-versus-implementation disputes,
covering who rules, within what period, what the possible dispositions are
(the specification is corrected, the library is corrected, or the construct is
declared implementation-defined and documented as such), and where the ruling
is recorded so the next implementer finds it. Issue \#6616 is the evidence that
the artifact alone is insufficient: the contradiction has been identified,
labelled, assigned, and milestoned, and it is still not adjudicated. An
ambiguity that is documented but unruled is not resolved.

The two halves reinforce each other. Without the procedure, the machine-readable
artifact records disagreements it cannot settle. Without the artifact, the
procedure has nothing precise to rule on and every dispute is re-litigated from
prose.

\noindent\textbf{Acceptance evidence.}
For a declared subset of the format --- not the whole format, which is a
multi-release scope --- the machine-readable description generates valid
fixtures that the reference library accepts and invalid fixtures it rejects,
and any divergence is either an identified library defect or an identified
description defect, with no third category. Every dispute raised in the
tracking period has a recorded disposition naming the authoritative artifact
and the change made, and the elapsed time from report to ruling is published.
A dispute closed by an implementation changing its output to match the
reference library, without a ruling on which artifact was authoritative, is
recorded as \emph{unresolved} rather than closed.

\subsubsection*{Recommendation IND-S1: Maintain a differential conformance and compatibility corpus}

Independent writers should validate generated files against normative on-disk
invariants and a matrix of current and supported older reference readers.
Reference-library maintainers should preserve independent-writer fixtures that
exercise specification-permitted alternatives, not only files emitted by
libhdf5 itself. The shared corpus should include both positive fixtures, which
must be accepted, and negative fixtures, which contain a single intentional
inconsistency and must be rejected predictably.

At minimum, the program should:

\begin{itemize}
    \item round-trip generated files through \code{h5dump}, \code{h5debug},
          h5py/libhdf5, and the independent implementation;
    \item test every supported chunk-layout and index variant, datatype width,
          filter pipeline, external-resource form, and library-version bound;
    \item preserve two separate regression cases: \#6409 for sufficient but
          nonminimal widths and \#6534 for contradictory element sizes;
    \item carry a third fixture class for \emph{contested} readings --- a
          construct whose conformance is genuinely undetermined, such as a
          group written without a group info message under \#6616 --- recorded
          with the competing interpretations and the expected result under
          each, so that a dispute has a test rather than an argument and so
          that the eventual ruling has something to close;
    \item report the writer, writer version, format bounds, reader, reader
          version, and expected accept/reject result for every fixture; and
    \item inventory already-produced files when a writer defect is discovered,
          provide a deterministic repair, regeneration, migration, or supported
          compatibility disposition, and preserve checksums and provenance when
          bytes change.
\end{itemize}

Acceptance requires the current supported-reader matrix to accept every valid
fixture, reject every intentionally invalid fixture with a controlled and
documented diagnostic, and demonstrate a reproducible repair, regeneration,
migration, or accountable supported-compatibility disposition for the complete
defined legacy PureHDF population, preserving checksums and provenance when
bytes change.

The machine-readable specification project \code{hdf5-pickles} is a natural
source for generating and checking these fixtures~\cite{hdf5-pickles}, and
under \code{FMT-S1} it is also the candidate normative artifact rather than
only a generator of test material.

\end{document}
